\chapter{The lessons}\label{ch:examples}

miniVivisection comes with 17 short script lessons, 12 Grasshopper lessons, the Markdown text example and an overview of all of
them. Each lesson has a bug in it that one feature finds or fixes, with instructions at the top of the file: \emph{The bug},
\emph{Try this}, \emph{Without miniVivisection}. New to debugging? Start with \file{00_Examples_Guide.pdf} in the
\file{examples} folder; it explains debugging for beginners and lists the lessons.

\section{Where they are}

\begin{description}
  \item[In the Script Editor:] The left panel's \menu{Examples\>miniVivisection} (\setting{Lessons in the Examples tab}).
  \item[On disk:] The \file{examples} folder next to the \file{.rhp}, the Grasshopper lessons in \file{examples\bs grasshopper}.
\end{description}

\paragraph{Your own copies.} The Examples tab opens your own copies, in \file{\%LOCALAPPDATA\%\bs miniVivisection\bs lessons},
and the Grasshopper lessons and their script files are there in \file{grasshopper}: save your fixes there, the shipped
lessons stay as they are. A copy you haven't changed follows a new version of the lesson. When a new miniVivisection
brings a new version of a lesson you changed, it is put beside your copy as \file{(name) (new)}, for example
\file{gh3_stop_a_hang (new).gh}. To start a lesson over, delete your copy; it is copied again when the editor next opens.
Open the Grasshopper copies with Grasshopper's \menu{File\>Open}.

\paragraph{How to do a lesson.}
\begin{enumerate}[label=\arabic*.]
  \item Open it, and read the text at the top (and a \emph{Needs:} line under the title when it needs something first).
  \item Set a breakpoint on the line marked \code{\# <-- BREAKPOINT}, and press \key{F5}, unless \emph{Try this} says
        otherwise: lessons 15 to 17 start with \ui{Debug} and no breakpoint; lesson 04 has two breakpoints.
  \item Follow the steps. When you have found the bug, stop, fix the code, and run again to check.
\end{enumerate}

\section{Where to find them}\label{sec:findexamples}

\begin{description}
  \item[The miniVivisection toolbar] in Rhino, shown when the plug-in is installed (\setting{Help, examples and settings toolbar}):
        \ui{Manual} -- the PDF; right-click: the Markdown in the Script Editor --, \ui{Examples} -- the overview of all
        examples, your copy of \file{00_overview.3dm} --, \ui{Examples in Grasshopper} -- Grasshopper started, its overview
        \file{gh0_overview.gh} opened --, \ui{Lessons folder} -- your copies, in Explorer --,
        \ui{Settings} -- Rhino's Options, at the miniVivisection page --, \ui{MarkdownText} and \ui{ProtectModel}.
  \item[The commands] \code{miniVivisectionManual} (\code{-miniVivisectionManual Markdown}: the Markdown),
        \code{miniVivisectionExamples}, \code{miniVivisectionGrasshopperExamples}, \code{miniVivisectionLessons},
        \code{miniVivisectionSettings} -- the toolbar's buttons, without it.
  \item[V°Examples, V°Manual, V°Settings] in Grasshopper, tab \menu{miniVivisection\>Help}, each a button (or a
        double-click): V°Examples opens \file{gh0_overview.gh} (your copy; shown when it is open), its menu the lessons
        folder; V°Manual the manual, its menu the Markdown in the Script Editor; V°Settings the settings
        (\setting{Help, examples and settings components} shows or hides all three; from Grasshopper's next start).
  \item[The Script Editor:] Its \ui{Examples} tab, folder \ui{miniVivisection}: the script lessons.
\end{description}

\section{The overview}

\file{00_overview.3dm} (Markdown text in Rhino) and \file{grasshopper\bs gh0_overview.gh} (a Panel) -- click a
link -- link to every lesson: the 17 script lessons open in the Script Editor, the Grasshopper lessons in Grasshopper, the
Markdown text example in Rhino. Each links to the other. Both are protected and say so in their text: \file{00_overview.3dm}
by \code{ProtectModel} -- move its V° logo, and its line says \emph{modified}; its link reverts --, \file{gh0_overview.gh} by a
V°Protect, with a slider to change.

\section{Script lessons}

\begin{xltabular}{\linewidth}{@{}P{0.37\linewidth}L>{\raggedleft\arraybackslash}p{0.07\linewidth}@{}}
\toprule
\textsf{\textbf{File}} & \textsf{\textbf{Feature}} & \\ \midrule
\lessonfile{01_watch_geometry}          & V°Watch: variables drawn in Rhino & \see{ch:watch} \\
\lessonfile{02_values_in_the_code}      & tooltips, line-end values, “was”, Tree; stepping & \see{ch:code} \\
\lessonfile{03_bake_before_a_crash}     & Bake & \see{sec:bake} \\
\lessonfile{04_edit_in_rhino}           & Edit in Rhino, Auto edit & \see{sec:editinrhino} \\
\lessonfile{05_conditional_breakpoints} & conditions, hit counts, log messages & \see{ch:breakpoints} \\
\lessonfile{06_watch_expressions}       & expressions in watches & \see{sec:addwatch} \\
\lessonfile{07_raw_data_interpret_as}   & Interpret as & \see{sec:interpretas} \\
\lessonfile{08_python_terminal}         & V°Python & \see{sec:vpython} \\
\lessonfile{09_step_over_other_files}   & stepping across files & \see{sec:debugging} \\
\lessonfile{10_visualizers}             & visualizers & \see{ch:visualizers} \\
\lessonfile{11_rhino_usable_while_paused} & FreePause & \see{sec:freepause} \\
\lessonfile{12_powershell_terminal}     & V°PowerShell & \see{sec:vpowershell} \\
\lessonfile{13_claude}                  & V°Claude & \see{sec:claude} \\
\lessonfile{14_threads_async}           & V°Threads, thread conditions, watches in every thread & \see{ch:threads} \\
\lessonfile{15_deadlock}                & Where is everyone?, Break now & \see{sec:vthreads} \\
\lessonfile{16_long_running_watch}      & Break now in a running script, fix it while paused, run on & \see{sec:breaknow} \\
\lessonfile{17_runtime_error}           & the debugger stops at a runtime error; values and call stack; Stop and record & \see{sec:runpythonscript} \\
\bottomrule
\end{xltabular}

\section{Grasshopper lessons}

Twelve definitions in \file{examples\bs grasshopper}, each with its lesson in a Panel on the canvas, ending with links to
the lesson before and after and to the overview, and under it a V°Protect: changed something? \ui{Revert} puts the lesson
back as it was shipped. GH~10's script file, \file{gh10_spiral.py}, is next to them.

\begin{xltabular}{\linewidth}{@{}P{0.37\linewidth}L>{\raggedleft\arraybackslash}p{0.07\linewidth}@{}}
\toprule
\textsf{\textbf{File}} & \textsf{\textbf{Feature}} & \\ \midrule
\lessonfile{gh1_which_branch_failed}    & failing runs, Copy failing branches, Debug failing run N & \see{sec:failingruns} \\
\lessonfile{gh2_the_one_wrong_item}     & Always debug, a breakpoint condition, the run badge, a slider moved while paused & \see{sec:runbadge} \\
\lessonfile{gh3_stop_a_hang}            & \keys{Ctrl+Alt+B} with the component's editor closed & \see{sec:ghbreak} \\
\lessonfile{gh4_v_edit}                 & V°Edit & \see{sec:vedit} \\
\lessonfile{gh5_runtime_error}          & the debugger stops at a runtime error in a script component (Always debug) & \see{sec:rewiring} \\
\lessonfile{gh6_threads_in_a_component} & threads in a component: results in finishing order, V°Threads, the freeze guard & \see{sec:ghasync} \\
\lessonfile{gh7_parallel_run}           & \code{ghpythonlib.parallel}: \code{parallel.run} that doesn't flatten, the pause at the error, a breakpoint inside the worker function & \see{sec:ghasync} \\
\lessonfile{gh8_long_running}           & a long-running component: \keys{Ctrl+Alt+B} pauses it (Always debug), a fix in V°Python, it runs on & \see{sec:ghbreak} \\
\lessonfile{gh9_script_from_a_panel}    & a script from a Panel: edited in the Script Editor, saved back into the Panel & \see{sec:scriptinput} \\
\lessonfile{gh10_script_from_a_file}    & a script file named in a Panel: \keys{Ctrl+Alt+B} stops the hang and opens the file at its line & \see{sec:scriptinput} \\
\lessonfile{gh11_markdown_notes}        & notes in Markdown, rendered in a Panel; editing the source where you double-click; a script Panel edited in place; links to components, a canvas view and the Rhino example & \see{sec:panelmarkdown} \\
\lessonfile{gh12_markdown_parser}       & Markdown with Rhino's parser: task boxes, nested lists and quotes, a footnote; a task and a table that need fixing & \see{sec:panelmarkdown} \\
\bottomrule
\end{xltabular}

\section{The Markdown text example}

In \lessonfile{00_overview.3dm}, right of the overview: Markdown made into Rhino text by \code{MarkdownText}. Double-click any of it to edit its
Markdown, and click its links: a named view, a viewport, the V° logo, GH~11's \emph{Ring} component and canvas
view, lesson 01, the overview, a web page (\cref{ch:markdowntext,ch:links}).
